\begin{center}\hypertarget{native-228-article165}{}\pdfbookmark[2]{165.}{article-165}165.\end{center}
Methodus simplicissima, orbitam proxime iam cognitam quatuor observationibus
adaptandi, haec esse videtur. Sint $x$, $y$ distantiae approximatae \SourceLeaf{229}corporis
coelestis a terra in duabus observationibus completis: harum adiumento
computentur loci respondentes heliocentrici, atque hinc ipsa elementa: ex his
dein elementis longitudines vel ascensiones rectae geocentricae pro duabus
reliquis observationibus. Quae si forte cum observatis conveniunt, elementa
nulla amplius correctione egebunt: sin minus, differentiae $X$, $Y$ notabuntur,
idemque calculus iterum bis repetetur, valoribus ipsarum $x$, $y$ paullulum mutatis.
Ita prodibunt tria systemata valorum quantitatum $x$, $y$ atque differentiarum
$X$, $Y$, unde per praecepta art. 120 valores correcti quantitatum $x$, $y$
eruentur, quibus valores $X=0$, $Y=0$ respondebunt. Calculo itaque simili huic
quarto systemati superstructo elementa emergent, per quae omnes quatuor
observationes rite repraesentabuntur.

Ceterum, siquidem eligendi potestas datur, eas observationes completas retinere
praestabit, e quibus situm orbitae maxima praecisione determinare licet, proin
duas observationes extremas, quoties motum heliocentricum $90$ graduum minoremve
complectuntur. Sin vero praecisione aequali non gaudent, earum latitudines vel
declinationes sepones, quas minus exactas esse suspicaberis.
